% Code listings
%\usepackage[scaled]{beramono}
\usepackage{fontspec}

\makeatletter
\def\verbatim@nolig@list{}
\makeatother

\setmonofont{JetBrains Mono}[Contextuals=Alternate]

\usepackage{ocr}
\usepackage[T1]{fontenc}
\usepackage[skins,breakable,listings]{tcolorbox}

\usepackage{accsupp}
\newcommand{\noncopyable}[1]{%
    \BeginAccSupp{method=escape,ActualText={}}%
    #1%
    \EndAccSupp{}%
}

\lstdefinelanguage{kotlin}{
    comment=[l]{//},
    commentstyle={\color{gray}\ttfamily},
    emph={delegate, filter, firstOrNull, forEach, it, lazy, mapNotNull, println, repeat, assert, with, head, tail, len, return@},
    numberstyle=\noncopyable,
    emphstyle={\color{OrangeRed}},
    identifierstyle=\color{black},
    keywords={abstract, actual, as, as?, break, by, class, companion, continue, data, do, dynamic, else, enum, expect, false, final, for, fun, get, if, import, in, infix, interface, internal, is, null, object, open, operator, override, package, private, public, return, sealed, set, super, suspend, this, throw, true, try, typealias, val, var, vararg, when, where, while, tailrec, reified},
    keywordstyle={\color{NavyBlue}\bfseries},
    morecomment=[s]{/*}{*/},
    morestring=[b]",
    morestring=[s]{"""*}{*"""},
    ndkeywords={@Deprecated, @JvmField, @JvmName, @JvmOverloads, @JvmStatic, @JvmSynthetic, Array, Byte, Double, Float, Boolean, Int, Integer, Iterable, Long, Runnable, Short, String},
    ndkeywordstyle={\color{BurntOrange}\bfseries},
    sensitive=true,
    stringstyle={\color{ForestGreen}\ttfamily},
    literate={`}{{\char0}}1
}

\lstdefinelanguage{docker}{
    keywords={FROM, RUN, COPY, ADD, ENTRYPOINT, CMD,  ENV, ARG, WORKDIR, EXPOSE, LABEL, USER, VOLUME, STOPSIGNAL, ONBUILD, MAINTAINER},
    keywordstyle=\color{blue}\bfseries,
    numberstyle=\noncopyable,
    identifierstyle=\color{black},
    sensitive=true,
    comment=[l]{\#},
    commentstyle=\color{purple}\ttfamily,
    stringstyle=\color{red}\ttfamily,
    morestring=[b]',
    morestring=[b]"
}

\lstdefinelanguage{BashPrompt}{%
    keywordstyle=\color{blue},
    basicstyle=\ttfamily,
    literate={\$}{{\bfseries\textcolor{blue}{\noncopyable\$}}}1
    {~}{{\bfseries\textcolor{blue}{\noncopyable\textasciitilde}}}1,
}

\lstdefinelanguage{DtsPrompt}{%
    keywordstyle=\color{blue},
    basicstyle=\ttfamily,
    literate={dt> }{{\bfseries\textcolor{blue}{\noncopyable{dt> }}}}1
}

\lstdefinelanguage{XML}{
    basicstyle=\ttfamily,
    numberstyle=\noncopyable,
    morestring=[s][\color{black}]{"}{"},
    morestring=[s][\color{black}]{>}{<},
    stringstyle=\color{black},
    identifierstyle=\color{blue}\bfseries,
    moreidentifier={launch, node},
    keywordstyle=\color{brown},
    morekeywords={name, pkg, type, output},
}

\lstdefinelanguage{bnf}{
    comment=[l]{//},
    commentstyle={\color{gray}\ttfamily},
    numberstyle=\noncopyable,
    basicstyle=\ttfamily,
    morestring=[s][\color{RoyalPurple}\bfseries]{"}{"},
    emphstyle={\color{ForestGreen}\bfseries},
    moreemph={TYPE, FIELD, SEPARATOR, CONSTANT, COMMENT},
    identifierstyle=\color{blue}\bfseries,
    keywordstyle={\color{NavyBlue}\bfseries},
    morekeywords={property, recover_property, rosInterfaceFile},
    ndkeywords={getType, get Key, getValue, getName, setName, getNameIdentifier},
    ndkeywordstyle={\color{black}\itshape},
}

%% Adjustments

\lstdefinestyle{barelisting}{
    language=kotlin,
    basicstyle=\ttfamily\scriptsize,
    numberstyle=\footnotesize\noncopyable,
    backgroundcolor=\color{slightgray},
    frame=single,
    tabsize=2,
    rulecolor=\color{black!30},
    breaklines=true,
    inputencoding=utf8
}

\usepackage{lipsum}
\tcbuselibrary{listings,breakable}

\newtcblisting{unbreakablekotlin}[1][]{%
    breakable=false,
    listing only,
    boxsep=-1.5pt,
    top=-1pt,
    bottom=-0.5pt,
    listing options={
        language=kotlin,
        basicstyle=\ttfamily\footnotesize,
        numberstyle=\footnotesize\noncopyable,
        tabsize=2,
        numbers=right,
        breaklines=true,
        inputencoding=utf8,
        escapeinside={(*@}{@*)},
        #1
    },
    underlay={
        \path[draw=none] (interior.north west) rectangle node[white]{\includegraphics[width=5mm]{../figures/kotlin_file.png}} ([xshift=-10mm,yshift=-10mm]interior.north west);
    }
}

\newtcblisting{kotlinlisting}[1][]{%
    listing options={
        language=kotlin,
        basicstyle=\ttfamily\tiny,
        numberstyle=\footnotesize\noncopyable,
        showstringspaces=false,
        tabsize=2,
        breaklines=true,
        numbers=right,
        inputencoding=utf8,
        escapeinside={(*@}{@*)},
        #1
    },
    underlay unbroken and first={%
        \path[draw=none] (interior.north west) rectangle node[white]{\includegraphics[width=4mm]{../figures/kotlin_file.png}} ([xshift=-10mm,yshift=-10mm]interior.north west);
    }
}

\newtcblisting{pclisting}[1][]{%
    listing options={
        language=BashPrompt,
        basicstyle=\ttfamily\footnotesize,
        numberstyle=\noncopyable,
        tabsize=2,
        numbers=right,
        inputencoding=utf8,
        escapeinside={(*@}{@*)},
        #1
    },
    underlay unbroken and first={%
        \path[draw=none] (interior.north west) rectangle node[white]{\includegraphics[width=5mm]{../figures/laptop_icon.png}} ([xshift=-10mm,yshift=-10mm]interior.north west);
    }
}

\newtcblisting{flexlisting}[1][]{%
    left=20pt,
    listing options={
        language=DtsPrompt,
        basicstyle=\ttfamily\footnotesize,
        numberstyle=\noncopyable,
        numbers=right,
        tabsize=2,
        inputencoding=utf8,
        #1
    }
}

\newtcblisting{dtslisting}[1][]{%
    left=20pt,
    listing options={
        language=DtsPrompt,
        basicstyle=\ttfamily\footnotesize,
        tabsize=2,
        inputencoding=utf8,
        #1
    },
    underlay unbroken and first={%
        \path[draw=none] (interior.north west) rectangle node[white]{\includegraphics[width=6mm]{../figures/duckietown.png}} ([xshift=-10mm,yshift=-10mm]interior.north west);
    }
}

% Imitate syntax error
\usepackage{ulem}
\makeatletter
\def\uwave{\bgroup \markoverwith{\lower3.5\p@\hbox{\sixly \textcolor{red}{\char58}}}\ULon}
\font\sixly=lasy6 % does not re-load if already loaded, so no memory problem.
\makeatother

\newtcblisting{launchlisting}[1][]{%
    #1,
    listing options={
        language=XML,
        basicstyle=\ttfamily\footnotesize,
        numberstyle=\footnotesize\noncopyable,
        tabsize=2,
        numbers=right,
        breaklines=true,
        inputencoding=utf8,
        escapeinside={(*@}{@*)}
    },
    fonttitle=\ttfamily\small,
    underlay unbroken and first={%
        \path[draw=none] (interior.north west) rectangle node[white]{\includegraphics[width=5mm]{../figures/launch_icon.png}} ([xshift=-10mm,yshift=-9mm]interior.north west);
    }
}

\newtcblisting{dockerlisting}[1][]{%
    listing options={
        language=docker,
        basicstyle=\ttfamily\footnotesize,
        numberstyle=\footnotesize\noncopyable,
        tabsize=2,
        numbers=right,
        breaklines=true,
        inputencoding=utf8,
        escapeinside={(*@}{@*)}
    },
    underlay unbroken and first={%
        \path[draw=none] (interior.north west) rectangle node[white]{\includegraphics[width=6mm]{../figures/docker_icon.jpg}} ([xshift=-10mm,yshift=-9mm]interior.north west);
    }
}

\newtcblisting{bnflisting}[1][]{%
    #1,
    listing options={
        language=bnf,
        basicstyle=\ttfamily\footnotesize,
        numberstyle=\noncopyable,
        breaklines=true,
        numbers=right,
        showstringspaces=false,
    },
    fonttitle=\ttfamily\small,
    underlay unbroken and first={%
        \path[draw=none] (interior.north west) rectangle node[white]{\includegraphics[width=5mm]{../figures/bnf_file.png}} ([xshift=-10mm,yshift=-10mm]interior.north west);
    }
}

\newtcblisting{pythonlisting}[1][]{%
    #1,
    listing options={
        language=Python,
        basicstyle=\ttfamily\footnotesize,
        numberstyle=\noncopyable,
        breaklines=true,
        numbers=right,
        showstringspaces=false,
        keywordstyle=\color{blue}\bfseries,
        escapeinside={(*@}{@*)}
    },
    fonttitle=\ttfamily\small,
    underlay unbroken and first={%
        \path[draw=none] (interior.north west) rectangle node[white]{\includegraphics[width=5mm]{../figures/python_icon.png}} ([xshift=-10mm,yshift=-10mm]interior.north west);
    }
}

\newtcblisting{rpilisting}[1][]{%
    listing options={
        language=BashPrompt,
        basicstyle=\ttfamily\footnotesize,
        numberstyle=\noncopyable,
        tabsize=2,
        inputencoding=utf8,
        #1
    },
    underlay unbroken and first={%
        \path[draw=none] (interior.north west) rectangle node[white]{\includegraphics[width=8mm]{../figures/raspi_icon.png}} ([xshift=-10mm,yshift=-10mm]interior.north west);
    }
}

\tcbset{
    enhanced jigsaw,
    breakable,
    listing only,
    boxsep=-1pt,
    top=-1pt,
    bottom=-0.5pt,
    right=-0.5pt,
    overlay first={
        \node[black!50] (S) at (frame.south) {\Large\ding{34}};
        \draw[dashed,black!50] (frame.south west) -- (S) -- (frame.south east);
    },
    overlay middle={
        \node[black!50] (S) at (frame.south) {\Large\ding{34}};
        \draw[dashed,black!50] (frame.south west) -- (S) -- (frame.south east);
        \node[black!50] (S) at (frame.north) {\Large\ding{34}};
        \draw[dashed,black!50] (frame.north west) -- (S) -- (frame.north east);
    },
    overlay last={
        \node[black!50] (S) at (frame.north) {\Large\ding{34}};
        \draw[dashed,black!50] (frame.north west) -- (S) -- (frame.north east);
    },
    before={\par\vspace{10pt}},
    after={\par\vspace{\parskip}\noindent}
}

\newcommand*{\inlineimg}[1]{%
    \raisebox{-.3\baselineskip}{%
        \includegraphics[
            height=\baselineskip,
            width=\baselineskip,
            keepaspectratio,
        ]{#1}%
    }%
}

\definecolor{slightgray}{rgb}{0.90, 0.90, 0.90}

\usepackage{soul}
\makeatletter
\def\SOUL@hlpreamble{%
    \setul{}{3.0ex}%
    \let\SOUL@stcolor\SOUL@hlcolor%
    \SOUL@stpreamble%
}
\makeatother

\newcommand{\inline}[1]{%
    \begingroup%
    \sethlcolor{slightgray}%
    \hl{\ttfamily\footnotesize #1}%
    \endgroup
}

\newcommand{\tinline}[1]{%
    \begingroup%
    \sethlcolor{slightgray}%
    \hl{\ttfamily\tiny #1}%
    \endgroup
}
